\chapter{Аналитическая часть}

В аналитической разделе будут рассматриваются основные подходы к визуализации огня и сформулированы решения, принимаемые при разработке программного обеспечения

\section{Постановка задачи визуализации и структура сцены}

Рассматривается трёхмерная сцена, в которой необходимо визуализировать один или несколько источников пламени. Каждый источник огня моделируется как независимая система частиц, эволюционирующих во времени под действием заданных правил. Помимо систем частиц, сцена включает внешний источник света и виртуальную камеру, определяющую точку зрения наблюдателя.

Состояние сцены на дискретном шаге времени можно описать следующими компонентами:
\begin{enumerate}
	\item множеством систем частиц, каждая из которых задаёт отдельный источник пламени;
	\item одним точечным источником внешнего света с настраиваемыми положением и интенсивностью;
	\item виртуальной камерой с настраиваемыми положением и параметрами проекции;
	\item прямоугольной областью вывода, соответствующей окну приложения.
\end{enumerate}

Задача визуализации состоит в том, чтобы по текущему состоянию систем частиц, источника света и камеры формировать последовательность двумерных кадров, создающих впечатление наблюдения за трёхмерным огнём. При этом должны обеспечиваться:
\begin{enumerate}
	\item визуально правдоподобное поведение пламени;
	\item возможность независимой настройки параметров каждой системы частиц (количество частиц, скорость, размер, цвет, характеристики эмиссии);
	\item управляемость параметров внешнего освещения и виртуальной камеры в реальном времени;
\end{enumerate}

\section{Метод систем частиц для моделирования пламени}

Метод систем частиц основан на представлении сложного по форме и динамике объекта набором элементарных частиц, каждая из которых имеет заданное состояние и подчиняется правилам эволюции.\cite{Reeves1983} Для моделирования огня такой подход позволяет воспроизводить форму факела, мерцание и затухание пламени без физически точного моделирования процессов горения.

\subsection{Параметры частицы и динамика во времени}

Состояние одной частицы удобно описывать вектором
\begin{equation}
	x(t) = \bigl( \mathbf{p}(t), \mathbf{v}(t), \tau(t), s(t), \alpha(t), \mathbf{C}(t) \bigr),
\end{equation}
где:
\begin{enumerate}
	\item $\mathbf{p}(t) = (x(t), y(t), z(t))$ -- положение частицы в мировых координатах;
	\item $\mathbf{v}(t)$ -- вектор скорости;
	\item $\tau(t) \in [0, 1]$ -- нормированная доля оставшейся жизни частицы;
	\item $s(t)$ -- линейный размер частицы;
	\item $\alpha(t) \in [0, 1]$ -- коэффициент прозрачности;
	\item $\mathbf{C}(t)$ -- цветовой вектор (например, в пространстве RGB).
\end{enumerate}

На каждом шаге моделирования выполняется дискретное обновление состояния по времени с шагом $\Delta t$:
\begin{equation}
	\mathbf{p}(t + \Delta t) = \mathbf{p}(t) + \mathbf{v}(t)\,\Delta t,
\end{equation}
\begin{equation}
	\mathbf{v}(t + \Delta t) = \mathbf{v}(t) + \mathbf{a}(t)\,\Delta t,
\end{equation}
где $\mathbf{a}(t)$ -- результирующее ускорение, включающее, например, подъёмную силу, действующую вверх, и псевдослучайные возмущения, имитирующие турбулентность \cite{DemCG}.

Нормированная жизнь частицы уменьшается по закону:
\begin{equation}
	\tau(t + \Delta t) = \max\bigl(0,\; \tau(t) - k_{\tau}\,\Delta t\bigr),
\end{equation}
где $k_{\tau} > 0$ -- параметр, определяющий среднюю продолжительность существования частицы. При достижении $\tau = 0$ частица считается мёртвой и подлежит удалению из системы.

Визуальные параметры задаются как функции от $\tau$:
\begin{equation}
	s(t) = s_{0}\, f_{s}\bigl(\tau(t)\bigr), \quad
	\alpha(t) = f_{\alpha}\bigl(\tau(t)\bigr), \quad
	\mathbf{C}(t) = \mathcal{C}\bigl(\tau(t)\bigr),
\end{equation}
где $s_{0}$ -- базовый размер частицы, а функции $f_{s}$, $f_{\alpha}$ и $\mathcal{C}$ задают законы изменения размера, прозрачности и цвета в зависимости от стадии ``жизни'' частицы.\cite{EckCG}

Для пламени функция $\mathcal{C}(\tau)$ обычно определяет переход от более горячих оттенков (белый, жёлтый) при больших значениях $\tau$ к более холодным (оранжевый, красный, дымный) при малых значениях $\tau$.

\subsection{Параметры системы частиц}

Система частиц задаёт множество частиц и параметры их рождения в сцене. В общем виде систему можно описать набором:
\begin{equation}
	P = \bigl( \Omega, \lambda, N_{\max}, \Theta \bigr),
\end{equation}
где:
\begin{enumerate}
	\item $\Omega \subset \mathbb{R}^{3}$ -- область эмиссии частиц (эмиттер), например круг или диск на плоскости;
	\item $\lambda$ -- средняя интенсивность эмиссии (число рождающихся частиц в единицу времени);
	\item $N_{\max}$ -- максимальное число частиц в системе;
	\item $\Theta$ -- набор параметров, определяющих начальные распределения скоростей, времени жизни и базовых цветов частиц.
\end{enumerate}

Процесс эволюции системы частиц на каждом шаге включает:
\begin{enumerate}
	\item эмиссию новых частиц в области $\Omega$ до достижения $N_{\max}$;
	\item обновление состояния всех живых частиц по формулам (2.2)–(2.4);
	\item пересчёт визуальных параметров $s(t)$, $\alpha(t)$ и $\mathbf{C}(t)$ по формулам (2.5);
	\item удаление частиц с $\tau = 0$.
\end{enumerate}

Такая схема позволяет описывать различные типы источников пламени, изменяя только параметры эмиттера и законы изменения визуальных характеристик.

\section{Параметры перспективной камеры}

Математически перспективная проекция задаётся усечённой пирамидой видимости (view frustum), которая определяется набором параметров \cite{DemCG}:
\begin{itemize}
	\item $\mathbf{p}_C = (x_c, y_c, z_c)$ -- положение камеры в мировой системе координат.
	\item $\mathbf{d}_C$ -- нормализованный вектор направления взгляда.
	\item $\mathbf{u}_C$ -- вектор «вверх», задающий ориентацию камеры вокруг оси взгляда.
	\item $\varphi$ -- угол поля зрения по вертикали . Определяет «широкоугольность» камеры. Большие значения $\varphi$ дают эффект панорамы, малые --- эффект телеобъектива.
	\item $a = w / h$ -- соотношение сторон целевого изображения.
	\item $n$, $f$ -- расстояния до ближней и дальней плоскостей отсечения. Объекты, находящиеся ближе $n$ или дальше $f$, не отображаются.
\end{itemize}

На основе этих параметров формируются матрицы преобразования: матрица вида $M_{view}$ (переход из мировых координат в координаты камеры) и матрица перспективной проекции $M_{proj}$.

\section{Альтернативные подходы к визуализации огня}

Системы частиц являются лишь одним из распространённых подходов к визуализации огня. Для полноты картины кратко рассмотрены альтернативы и обоснован выбор именно данного метода.

\subsection{Текстурные и спрайтовые анимации}

Текстурные и спрайтовые анимации основаны на использовании заранее подготовленных последовательностей изображений огня, отображаемых на плоских полигонах или билбордах \cite{FoleyVanDam}. Движение пламени формируется за счёт смены кадров анимации.

Преимущества:
\begin{enumerate}
	\item крайне простая реализация;
	\item минимальная вычислительная сложность;
	\item предсказуемое визуальное качество.
\end{enumerate}

Недостатки:
\begin{enumerate}
	\item привязка к заранее подготовленным анимациям;
	\item ограниченные возможности настройки формы и поведения пламени;
	\item слабая интеграция с динамическим освещением.
\end{enumerate}

\subsection{Объёмные методы и объёмный рендеринг}

В объёмных методах огонь представляется в виде трёхмерного поля плотности, температуры или коэффициентов эмиссии и поглощения, заданного на регулярной воксельной сетке \cite{VolumeRendering}. Визуализация выполняется методом объёмного рендеринга, предполагающим интегрирование вклада среды вдоль луча зрения.

Преимущества:
\begin{enumerate}
	\item высокая степень реализма, возможность учёта самозатенения и рассеяния;
	\item естественное представление сложной структуры пламени.
\end{enumerate}

Недостатки:
\begin{enumerate}
	\item необходимость хранения и обновления трёхмерных полей;
	\item значительная вычислительная сложность процедуры рендеринга;
	\item повышенные требования к аппаратным ресурсам.
\end{enumerate}

\subsection{Физически обоснованное моделирование на основе уравнений гидродинамики}

Физически обоснованные модели описывают движение и нагрев газовой среды с помощью уравнений гидродинамики и уравнений переноса энергии, которые затем используются для визуализации пламени \cite{FoleyVanDam}.

Преимущества:
\begin{enumerate}
	\item физическая достоверность поведения огня;
	\item возможность моделирования взаимодействия потока с препятствиями.
\end{enumerate}

Недостатки:
\begin{enumerate}
	\item сложность реализации и настройки численных методов;
	\item высокие требования к вычислительным ресурсам;
	\item неочевидная связь физических параметров с удобными для пользователя визуальными параметрами.
\end{enumerate}

\section{Модели освещения}

Модель освещения определяет, как параметры источников света и визуальные свойства объектов сцены преобразуются в наблюдаемый цвет на экране. Для огня важно учесть, что он является самосветящейся полупрозрачной средой, на которую дополнительно влияет внешний источник света.

\subsection{Классические локальные модели освещения}

В традиционных трёхмерных сценах широко используются локальные модели освещения, такие как диффузная модель Ламберта и модель Фонга.\cite{FoleyVanDam} В диффузной модели яркость точки поверхности определяется выражением
\begin{equation}
	I_{\text{diff}} = I_{L} k_{d} \max\bigl(0, \mathbf{n} \cdot \mathbf{l}\bigr),
\end{equation}
где $I_{L}$ -- интенсивность источника света, $k_{d}$ -- диффузный коэффициент материала, $\mathbf{n}$ -- единичный вектор нормали к поверхности, $\mathbf{l}$ -- единичный вектор, направленный на источник.

Модель Фонга добавляет зеркальную составляющую:
\begin{equation}
	I_{\text{spec}} = I_{L} k_{s} \max\bigl(0, \mathbf{r} \cdot \mathbf{v}\bigr)^{n},
\end{equation}
где $k_{s}$ -- зеркальный коэффициент, $\mathbf{r}$ -- направление отражённого луча, $\mathbf{v}$ -- направление на наблюдателя, $n$ -- показатель блеска.\cite{FoleyVanDam}

Эти модели предполагают наличие определённой геометрической поверхности и нормали, что не вполне соответствует частицам как элементам объёмной полупрозрачной среды.

\subsection{Освещение самосветящейся полупрозрачной среды}

Огонь можно рассматривать как самосветящуюся среду, в которой каждая частица сама испускает свет, а внешний источник света лишь модифицирует наблюдаемую яркость \cite{VolumeRendering}. В приближении цвет частицы определяется эмиссионной составляющей и коэффициентом, зависящим от внешнего освещения.

Пусть $\mathbf{C}_{\text{emiss}}$ -- эмиссионный цвет частицы, определяемый её временем жизни и базовым цветом системы, а внешнее освещение учитывается через коэффициент затухания по расстоянию до собственного источника света:
\begin{equation}
	a(d) = \frac{1}{1 + k d^{2}},
\end{equation}
где $d$ --- расстояние от частицы до источника, $k > 0$ -- параметр затухания.\cite{FoleyVanDam}

Тогда итоговый цвет частицы можно задать как
\begin{equation}
	\mathbf{C}_{\text{final}} =
	\mathbf{C}_{\text{emiss}} \odot \bigl( \mathbf{1} + \beta\, a(d) \bigr),
\end{equation}
где $\mathbf{1} = (1, 1, 1)$, $\beta$ -- коэффициент влияния внешнего света, а символ $\odot$ обозначает поэлементное умножение компонент цвета.\cite{EckCG}

\section{Алгоритмы удаления невидимых поверхностей}

При визуализации трёхмерной сцены необходимо определить, какие элементы должны быть видимы в каждом пикселе изображения. Это приводит к задаче удаления невидимых линий и поверхностей \cite{DemCG}. В рассматриваемой сцене присутствуют как потенциально непрозрачные элементы (фон), так и большое количество полупрозрачных частиц, моделирующих пламя.

\subsection{Постановка задачи видимости}

После перехода к координатам камеры и применения проекции каждая вершина полигона или частицы имеет экранные координаты $(x, y)$ и значение глубины $z$, характеризующее расстояние до наблюдателя \cite{RogersAdams}. Для каждого пикселя $(x, y)$ требуется выбрать фрагмент с минимальным значением $z$ среди всех фрагментов, проецирующихся в данный пиксель:
\begin{equation}
	F^{\ast}(x, y) = \arg\min_{F_{k}} z_{k}(x, y).
\end{equation}

Если грань объекта лежит на плоскости
\begin{equation}
	A x + B y + C z + D = 0,
\end{equation}
то при фиксированных $x$ и $y$ глубина точки этой грани может быть получена как
\begin{equation}
	z(x, y) = -\frac{A x + B y + D}{C},
\end{equation}
что используется в ряде алгоритмов определения видимости.\cite{RogersAdams}

\subsection{Алгоритм Робертса}

Алгоритм Робертса относится к числу ранних методов удаления невидимых линий и поверхностей и основан на попарном анализе взаимного перекрытия полигонов в проекции \cite{RogersAdams}. В общем виде алгоритм включает:
\begin{enumerate}
	\item вычисление уравнений плоскостей для всех граней объектов сцены;
	\item проецирование граней на плоскость экрана и определение областей пересечения их проекций;
	\item сравнение значений глубины $z_{i}(x, y)$ и $z_{j}(x, y)$ для граней $P_{i}$ и $P_{j}$ в характерных точках областей пересечения;
	\item при необходимости -- разбиение граней на части, соответствующие видимым и невидимым областям.
\end{enumerate}

Математически видимость в точке экрана определяется сравнением значений
\begin{equation}
	z_{i}(x, y) = -\frac{A_{i} x + B_{i} y + D_{i}}{C_{i}}, \quad
	z_{j}(x, y) = -\frac{A_{j} x + B_{j} y + D_{j}}{C_{j}}.
\end{equation}

Преимущество алгоритма Робертса -- учёт взаимного перекрытия граней; недостатки связаны с вычислительной сложностью и трудоёмкостью реализации при большом количестве полигонов \cite{RogersAdams}.

\subsection{Алгоритм Варнока}

Алгоритм Варнока относится к подрастровым методам и основан на рекурсивном разбиении области экрана на более мелкие прямоугольники.\cite{FoleyVanDam} Для исходного окна изображения $R_{0}$ определяется множество полигонов, проецирующихся в эту область. Далее выполняются следующие шаги:
\begin{enumerate}
	\item если в $R_{0}$ не попадает ни один полигон, область заполняется фоновым цветом;
	\item если в $R_{0}$ проецируется один полигон и полностью покрывает его, область заполняется цветом этого полигона;
	\item в общем случае $R_{0}$ разбивается, как правило, на четыре подокна, и процедура рекурсивно повторяется для каждого подокна.
\end{enumerate}

При этом анализируется пересечение полигонов с подокнами, а глубина в характерных точках (например, в центре подокна) вычисляется по уравнению вида \cite{RogersAdams}.

Алгоритм Варнока обеспечивает адаптивное уточнение решения в областях с высокой геометрической сложностью, но требует значительного числа операций пересечения и обрезки полигонов, что затрудняет его применение в динамических сценах с большим числом примитивов.

\subsection{Алгоритм z-буфера}

Алгоритм z-буфера реализует решение задачи видимости непосредственно в растровом пространстве \cite{RogersAdams}. Для каждого пикселя $(x, y)$ поддерживаются значения глубины ближайшего уже отрисованного фрагмента $z_{\text{buf}}(x, y)$ и его цвета $C_{\text{buf}}(x, y)$.

Начальная инициализация имеет вид
\begin{equation}
	z_{\text{buf}}(x, y) = z_{\max}, \quad
	C_{\text{buf}}(x, y) = C_{\text{background}},
\end{equation}
где $z_{\max}$ -- значение глубины на дальней плоскости отсечения.

При обработке нового фрагмента с глубиной $z_{\text{new}}$ и цветом $C_{\text{new}}$ выполняется сравнение:
\begin{equation}
	\text{если } z_{\text{new}} < z_{\text{buf}}(x, y),
\end{equation}
\begin{equation}
	\text{то } z_{\text{buf}}(x, y) := z_{\text{new}}, \quad
	C_{\text{buf}}(x, y) := C_{\text{new}},
\end{equation}
иначе фрагмент считается невидимым и отбрасывается.

Преимущества z-буфера:
\begin{enumerate}
	\item отсутствие необходимости глобальной сортировки объектов по глубине;
	\item масштабируемость при большом числе простых примитивов.\cite{RogersAdams}
\end{enumerate}

\section{Сравнительный анализ и выбор методов}
\label{subsec:comparison_and_choice}

Для обоснованного выбора ключевых методов визуализации был проведён сравнительный анализ по сформулированным критериям, наиболее значимым для задач моделирования огня системами частиц в реальном времени.

В таблице~\ref{tab:lighting_comparison} представлено сравнение рассмотренных моделей освещения.

\begin{table}[H]
	\centering
	\small
	\caption{Сравнительный анализ моделей освещения для систем частиц}
	\begin{tabular}{|p{4.5cm}|c|c|c|}
		\hline
		\textbf{Критерий / Модель} & \textbf{Локальные} & \textbf{Объёмная} & \textbf{Эмиссионная} \\
		\hline
		Адекватность физике огня & -- & + & + \\
		\hline
		Учёт самосвечения & -- & + & + \\
		\hline
		Учёт внешнего освещения & + & + & + \\
		\hline
		Вычислительная эффективность & + & -- & + \\
		\hline
	\end{tabular}
	\label{tab:lighting_comparison}
\end{table}

Классические локальные модели неприменимы к частицам из-за отсутствия у них поверхности и нормали. Объёмная модель физически наиболее адекватна, но обладает неприемлемо высокой вычислительной сложностью. Принятая эмиссионная модель обеспечивает баланс: она учитывает самосвечение и внешний свет, оставаясь при этом вычислительно эффективной, что соответствует требованиям интерактивного приложения.

В таблице~\ref{tab:hidden_surface_comparison} представлено сравнение рассмотренных алгоритмов удаления невидимых поверхностей.

\begin{table}[H]
	\centering
	\small
	\caption{Сравнительный анализ алгоритмов удаления невидимых поверхностей}
	\begin{tabular}{|p{4.5cm}|c|c|c|}
		\hline
		\textbf{Критерий / Алгоритм} & \textbf{Робертса} & \textbf{Варнока} & \textbf{Z-буфер (принятый)} \\
		\hline
		Производительность при большом $N$ & -- & -- & + \\
		\hline
		Поддержка полупрозрачности & -- & -- & + (с аддитивным смешиванием) \\
		\hline
		Применимость к динамической сцене & -- & -- & + \\
		\hline
	\end{tabular}
	\label{tab:hidden_surface_comparison}
\end{table}

Алгоритмы Робертса и Варнока демонстрируют низкую производительность при большом количестве примитивов и высокую сложность реализации, что делает их неподходящими для динамических систем частиц. Алгоритм z-буфера лишён этих недостатков: его производительность линейно зависит от числа примитивов, он адаптируется к динамическим сценам. В сочетании с аддитивным смешиванием цветов он обеспечивает корректное отображение полупрозрачных частиц пламени \cite{VolumeRendering}.

\section{Вывод}

В аналитической части были рассмотрены основные подходы к визуализации огня и сформулированы решения, принимаемые при разработке программного обеспечения.

\clearpage
